% chapters/21_multi_agent.tex
% Part of: AI / LLM / Agents / Hardware Notes

\section{Multi-agent Systems \& Orchestration}

\subsection{Orchestrator-subagent patterns}

In an orchestrator-subagent architecture, a primary (orchestrator) agent decomposes a
task and delegates sub-tasks to specialised child agents. The orchestrator maintains
the high-level plan; subagents handle isolated, bounded work.

\textbf{Claude Code / Cowork implementation} (observed directly):

\begin{itemize}
  \item The orchestrator has access to an \texttt{Agent} tool that launches a subagent
    as a subprocess. The subagent type determines which tools it receives
    (e.g.\ \texttt{general-purpose}, \texttt{Explore}, \texttt{Plan}).

  \item Subagents are \emph{isolated}: each receives a fresh context containing only
    the task prompt passed by the orchestrator, plus any tool results it generates.
    They do not share the orchestrator's conversation history.

  \item Subagents can be run \textbf{in parallel} (multiple \texttt{Agent} tool calls
    in one assistant turn) when sub-tasks are independent, dramatically reducing
    wall-clock time for multi-step research or multi-file operations.

  \item A subagent can be \textbf{resumed} by passing its \texttt{agent\_id} back to
    the \texttt{Agent} tool, preserving its full transcript for follow-up work without
    re-incurring input token costs.

  \item \textbf{Worktree isolation} (\texttt{isolation: "worktree"}) gives a subagent
    its own git worktree so it can make changes without touching the orchestrator's
    working tree; the worktree is auto-cleaned if no changes are made.
\end{itemize}

The key design tension is \emph{isolation vs.\ context sharing}: fully isolated
subagents avoid context pollution but may duplicate work; shared-context agents
collaborate more fluidly but risk token bloat and unintended interference.

\subsection{AutoGen, CrewAI, LangGraph \& MetaGPT}

\textbf{AutoGen} (Microsoft) builds agents around \emph{conversational turns}. Two
agents (or a human proxy and an agent) exchange messages to jointly solve a task.
\texttt{GroupChat} enables $n$-way conversations with a group manager that routes
each turn to the right participant. Execution is controlled via reply functions and
termination conditions.

\textbf{CrewAI} models a multi-agent system as a \emph{crew} of workers with defined
roles, backstories, goals, and tool access. Tasks are assigned to crew members and
can run sequentially or in parallel. Encourages role-based thinking (researcher,
writer, reviewer) rather than low-level code primitives.

\textbf{LangGraph} provides a \emph{stateful graph} execution model. Agents and
tools are nodes; edges encode control flow (sequential, conditional, cyclic):
\begin{itemize}
  \item State is a typed dict threaded through every node, making data flow
    inspectable and making bugs reproducible.
  \item Cycles are first-class: ReAct-style Thought/Action loops map naturally to
    a graph edge pointing back to an earlier node.
  \item \textbf{Human-in-the-loop:} interrupt nodes pause execution for human
    review/approval before the graph continues.
  \item Integrates with MCP servers via \texttt{langchain-mcp-adapters} for
    plug-and-play tool server compatibility.
\end{itemize}

\textbf{MetaGPT} simulates a \emph{virtual software company} using SOPs:
\begin{itemize}
  \item Roles: product manager, architect, project manager, engineer, QA engineer.
  \item Input: one-line requirement. Output: PRD, architecture doc, API specs, code,
    tests.
  \item \textbf{Standard Operating Procedures (SOPs):} human workflows encoded as
    prompt sequences; each agent's output is the structured input to the next.
  \item \textbf{Publish-subscribe message pool:} all inter-agent messages are
    structured artefacts (documents, diagrams) stored in a global pool, decoupling
    senders from receivers and reducing natural-language parsing errors.
\end{itemize}

\begin{center}
\begin{tikzpicture}[
  role/.style={rectangle, rounded corners=3pt, draw=aiblue, fill=infobg,
               minimum width=1.8cm, minimum height=0.7cm, font=\scriptsize\bfseries,
               align=center},
  arr/.style={-{Stealth[length=5pt]}, thick, color=aigray}
]
  \node[role] (pm)   {Product\\Manager};
  \node[role, right=1.3cm of pm]   (arch)  {Architect};
  \node[role, right=1.3cm of arch] (projm) {Project\\Mgr};
  \node[role, right=1.3cm of projm](eng)   {Engineer};
  \node[role, right=1.3cm of eng]  (qa)    {QA};
  \draw[arr] (pm)   -- node[above,font=\tiny]{PRD}    (arch);
  \draw[arr] (arch) -- node[above,font=\tiny]{Design} (projm);
  \draw[arr] (projm)-- node[above,font=\tiny]{Tasks}  (eng);
  \draw[arr] (eng)  -- node[above,font=\tiny]{Code}   (qa);
  \node[above=0.85cm of projm, font=\small\itshape, color=aigray]
    {MetaGPT: one-line requirement $\rightarrow$ full codebase};
\end{tikzpicture}
\end{center}

\textbf{MAS network topologies:}
\begin{itemize}
  \item \textbf{Centralized:} a central coordinator holds global knowledge and routes
    tasks to agents. Easy to monitor; single point of failure.
  \item \textbf{Decentralized:} agents share state only with their neighbours.
    More fault-tolerant and modular; coordination is harder.
  \item \textbf{Hierarchical:} tree structure with varying autonomy per layer.
    Natural for enterprise workflows where high-level strategy delegates to
    operational agents.
  \item \textbf{Holonic:} agents are grouped into \emph{holons} (a holon is a unit
    that is both a whole and a part). A holon's leading agent appears to the outside
    world as a single agent; internally it manages a sub-team.
  \item \textbf{Coalition / Team:} agents temporarily unite to achieve a shared goal,
    disbanding afterward. Teams are more tightly coupled; coalitions are ad hoc.
\end{itemize}

\subsection{Inter-agent communication}

Direct agent-to-agent integration suffers from an $n(n{-}1)/2$ integration-point
problem: $n$ agents each needing a bespoke connector to every other agent. Protocol
standards collapse this to $n$ connections to a shared message bus.

\textbf{Model Context Protocol (MCP)} --- Anthropic, 2024:
\begin{itemize}
  \item Focus: one model, many tools. Standardises how a single agent enriches its
    context via external servers.
  \item Transport: JSON-RPC over stdio or HTTP/SSE.
  \item Servers expose \emph{tools}, \emph{resources}, and \emph{prompts}.
  \item Limitation for multi-agent: no native agent-to-agent discovery, no delta
    streaming, no shared memory primitive.
\end{itemize}

\textbf{Agent Communication Protocol (ACP)} --- IBM (merged with A2A):
\begin{itemize}
  \item Focus: agent-to-agent communication across frameworks and organisations.
  \item Transport: REST over standard HTTP (no SDK required; cURL/browser compatible).
  \item \textbf{Async-first:} designed for long-running tasks; synchronous mode also
    supported.
  \item \textbf{Offline discovery:} agents embed metadata in their distribution
    packages so other agents can discover capabilities without a central registry.
  \item Enables orchestration across heterogeneous frameworks without custom glue.
\end{itemize}

\textbf{Agent2Agent Protocol (A2A)} --- Google:
\begin{itemize}
  \item Designed for cross-system and cross-organisation agent communication.
  \item More natural for Google-ecosystem integrations.
  \item Can coexist with ACP; each serves different primary use cases.
\end{itemize}

\begin{notebox}[title=MCP vs ACP in one line]
  MCP = one model $\leftrightarrow$ many tools. \quad
  ACP = many agents $\leftrightarrow$ many agents.
\end{notebox}

\subsection{Routers, state machines \& workflow engines}

As agent pipelines grow, structured control-flow primitives prevent LLM-driven
decision-making from becoming a runaway cost and reliability problem.

\textbf{Router:} decides which tool, subagent, or branch to invoke. Implementations:
\begin{itemize}
  \item LLM classifier --- flexible, expensive; use for open-ended routing.
  \item Embedding similarity match --- fast, moderate generality.
  \item Deterministic rule --- fastest, requires explicit case enumeration.
\end{itemize}

\textbf{State machine / workflow engine:} agent progression through phases is
explicit. Each state = a task phase (planning, retrieving, synthesising, reviewing);
transitions are triggered by events or conditions. LangGraph implements this directly.

Benefits over fully open-ended ReAct:
\begin{itemize}
  \item Deterministic paths for known task types (no unnecessary LLM calls).
  \item Explicit human-in-the-loop checkpoints at state transitions.
  \item Each state is independently testable.
  \item Defined failure handling per state, rather than global ``if error, retry.''
\end{itemize}

\begin{notebox}[title=Agent vs workflow]
  A \emph{workflow} uses fixed steps and deterministic transitions.
  An \emph{agent} makes dynamic decisions at each step.
  Many production ``agents'' are really workflows with one or two LLM-decision
  nodes --- often the better engineering choice.
\end{notebox}

\subsection{Swarm intelligence analogies}

Swarm intelligence studies how simple local rules among agents produce complex,
emergent global behaviour without central control. Canonical examples: ant colony
optimisation (ACO), particle swarm optimisation (PSO), bee colony algorithms.

Analogies to LLM multi-agent systems:
\begin{itemize}
  \item \textbf{Stigmergy:} ants communicate indirectly via pheromone trails left in
    the environment. In MAS, agents communicate indirectly via a shared message pool
    (MetaGPT), a vector database, or a shared workspace --- they don't call each other
    directly.
  \item \textbf{Emergent division of labour:} no agent is told to specialise; roles
    emerge from local interactions. In LLM swarms (e.g.\ Society of Mind prompting),
    assigning diverse ``perspectives'' to agents produces richer collective outputs
    than a single agent.
  \item \textbf{Robustness through redundancy:} decentralised MAS topologies mirror
    swarms: if one node fails, the system degrades gracefully rather than collapsing
    (contrast with centralised architectures with a single point of failure).
  \item \textbf{Limits of the analogy:} LLM agents are expensive, stateful, and
    context-limited; true swarms use thousands of cheap stateless units. Current
    multi-agent LLM systems typically operate at $n < 20$ agents.
\end{itemize}

